% =============================================================================
\section{Introduction}
\label{sec:intro}
% =============================================================================
\cgal{} (Computational Geometry Algorithm Library) is an open source
software project that provides access to efficient and reliable
implementations of geometric algorithms in the form of a C++
library~\cite{cgal:eb-21a}. \cgal{} has evolved through the years and
now represents the state of the art in computational geometry software. 
\cgal{} is used in a diverse range of domains requiring geometric 
computation such as computer graphics, scientific visualization, 
computer aided design and modeling, geographic information systems,
molecular biology, medical imaging, and many more. \cgal{} provides
a large number of components that cover a wide range of areas. The 
bindings for most components are supported now. The principles for the 
bindings described in this paper applies to all packages of \cgal{}
(and perhaps to other generic C++ libraries as well), and can be easily
adopted to support the bindings of the few components that are still 
missing.

There is a large community of users of \cgal{} 
packages in academia and industry, and these packages seem to be
helpful in many diverse projects. In all the projects that we are
aware of, the packages have been used by C++ experts or by users who
have been strongly supported by such experts from the \cgal{}
developer community, including
GeometryFactory.\footnote{GeometryFactory is the company established
  to market and support the usage of \cgal{} packages.} In our
experience, incorporating \cgal{} in education or in academic projects
lacking C++ experts has been rugged. We attribute the difficulties
that arise in such settings to the highly sophisticated usage of C++
in \cgal. The binding project that we describe in this paper is meant
to resolve these difficulties and make all \cgal{} packages more
accessible to a wider audience of programmers. We already see, and
report below on, first signs that the project is on a good trajectory
toward achieving this goal.

The software modules of \cgal{} rigorously adhere to the
generic programming paradigm---a discipline that consists of the
gradual lifting of concrete algorithms abstracting over details, while
retaining the algorithm semantics and efficiency. The software modules
of \cgal{} also follow the exact geometric-computation paradigm, which
simply amounts to ensuring that errors in predicate evaluations do not
occur; it guarantees robustness of the applied algorithms.  As a
result, the software of \cgal{} is robust despite the use of (inexact)
floating point arithmetic, found in standard hardware, it is complete
as it handles all degenerate cases, which typically are in abundance,
and it is efficient---all at the same time.

Generic programming identifies an abstraction that consists of a
formal hierarchy of polymorphic abstract requirements on data types,
referred to as \emph{concepts}, and a set of classes that conform
precisely to the specified requirements, referred to as
\emph{models}~\cite{a-gps-99}. A hierarchy of related concepts can be
viewed as a directed acyclic graph, where a node of the graph
represents a concept and an arc represents a refinement relation. An
arc directed from concept A to concept B indicates that concept B
refines concept A. When a class template is instantiated, each one of
its template parameters is substituted with a model of one or more
concepts associated with the template parameter.

% -----------------------------------------------------------------------------
\subsection{C++ Python Bindings}
\label{ssec:intro:bind}
% -----------------------------------------------------------------------------
Python, like many other languages, was originally defined by a
reference implementation. The original reference implementation of the
Python programming language is CPython. To date, CPython is still the
default and commonly used reference implementation of the language,
although new rivals are emerging, see, e.g., PyPy.\footnote{PyPy is a
  replacement for CPython; see \url{https://www.pypy.org}.} As the
Python language evolved, its definition has tightened and a well
defined specification has materialized.

Bindings are essentially wrapper libraries that bridge two programming
languages, so that a software module developed in one language can be
used in code written in another language, exploiting (i) the strengths
of both languages, and (ii) the availability of modules developed in
the former language. C++ Python bindings enable (i) the invocation of
C++ functions, and (ii)~the access to C++ variables from Python code,
taking advantage of Python's quick development cycles and C++'s high
performance.

Python is sufficiently efficient for many tasks; nevertheless
low-level code written in Python tends to be too slow, largely due to
Python's extremely dynamic nature. In particular, low-level
computational loops are simply infeasible. In addition, the magnitude
of existing and well-tested code in Fortran, C, or C++, is
enormous. The Python/C API, the module that enables the usage of
external modules developed in Fortran, C, or C++, from Python code, or
embedding Python code in code written in other languages, is very
rich.  Indeed, the Python/C API module of CPython exposes a vast
amount of details of CPython internals. It enables the efficient
exploitation of existing code in Fortran, C, or C++, as well as the
replacement of critical sections written in Python when speed is
essential. Wrapping existing code has traditionally been the domain of
Python experts due to the steep learning curve, which characterizes
its Python/C API. Although using such wrappers is possible without
ever knowing their internals, providing such wrappers creates a sharp
line between developers using Python and developers using C or C++
with the Python/C API. Several tools that facilitate the creation of
such wrappers have been introduced; A sample of such tools are listed
In Section~\ref{sec:comp}.

When running Python code that uses bindings for some C++ modules, one
or more libraries that provide the bindings must be accessible.  A
software module in C++ that adheres to the generic programming
paradigm consists of function and class templates; these templates are
instantiated at compile time of the binding libraries. In other words,
the types of the C++ objects that are bound, that is, instances
(instantiated types) of C++ function and class templates, must be
known when the bindings are generated.

% -----------------------------------------------------------------------------
\subsection{\cmake}
\label{ssec:intro:cmake}
% -----------------------------------------------------------------------------
\cmake{}~\cite{mh-mc-08} is an open-source, cross-platform suite of
tools designed to build, test, and package software. The suite of
\cmake{} tools were created to address the need for cross-platform
build environments for open-source projects. However, since it was
conceived over two decades ago, it has grown and become powerful,
rich, and flexible. It is now considered a modern tool; it is widely
spread and used by major commercial software products for their build
and test environments.

\cmake{}, for the most part, is a cross-platform build system
generator.  It is used to control the software build process using
simple platform and compiler independent configuration files, and it
generates native build environments of the user choice. Users build a
project by using \cmake{} to generate a build system for native tools
on their platforms. \cmake{} supports an interpreted, functional,
scripting language. The build process of a project is specified with
platform-independent \cmake{} text files, named
\myLstinline{CMakeLists.txt}, written in the \cmake{} language,
included in each directory of the source tree of the project. For
example, once the \cmake{} configuration files are in place for a
certain project, \cmake{} can be used to generate
\myLstinline{Makefile} files required by the \myLstinline{make}
utility of the GNU suite to compile a project on a Unix system using
\myLstinline{g++}. The command-line interface of the cross-platform
build system generator of \cmake{} is an executable called
\myLstinline{cmake}. It has a counterpart called
\myLstinline{cmake-gui}, which includes a graphical user interface
(GUI). To generate native project build files, the user invokes
\myLstinline{cmake} in a terminal and specifies the directory that
contains the root \myLstinline{CMakeLists.txt} file.

A key feature of \cmake{} is the ability to (optionally) place
compiler outputs (such as object files) outside the source tree. This
enables multiple and concurrent builds from the same source tree and
cross-compilation. It also keeps the source tree clean and ensures
that removing a build directory will not remove the source files.  We
exploit this feature, and the scripting capabilities to conveniently
generate the C++ Python bindings for several \cgal{} objects
concurrently.

% -----------------------------------------------------------------------------
\subsection{Bindings for \cgal}
\label{ssec:intro:cgal}
% -----------------------------------------------------------------------------
Instantiated types in \cgal{} are characterized by long instantiation
chains of C++ class templates. An instantiated type is a template that
has one or more template parameters and every parameter is substituted
by another type that is typically an instance of another
template. While the number of models of most concepts is small, the
number of potential types of objects that could be bound, and thus
must be supported by the bindings, is enormous. Offering bindings for
all these types in advance is practically impossible. Moreover, in
some cases models need to be extended with types provided by the
user. Naturally, bindings for user-extended models must be generated
dynamically.

\cgal{} uses \cmake~\cite{mh-mc-08} for the build process.  Since
version~5.0, \cgal{} is header-only by default, which means that an
application that depends on \cgal{} only depends on (i) the source
code of \cgal{} that resides in header files, (the names of which
typically end with the suffix \myLstinline{.hpp}), and (ii) native
configuration files generated by \cmake{} (in particular, the
execution of \myLstinline{cmake}). In other words, there is no need to
compile any library before an application that depends on \cgal{} and
written in C++ is built.\footnote{Some dependencies of \cgal{} might
  need to be installed in advance.} On the other hand, when running an
application written in Python that depends on \cgal{} the libraries
that contain the bindings must be accessible. Typically, those
bindings are generated ahead of time with little knowledge about the
application itself. This would impose a severe limitation in cases
like ours, where the number of objects to be bound is large. In cases
where the types of the objects to be bound are closely tied with the
application itself, generating bindings ahead of time is, naturally,
impossible. Our system enables easy and convenient (re)generation of
bindings, thus alleviating the burden caused by these problems.

% -----------------------------------------------------------------------------
\subsection{Contribution}
\label{ssec:intro:results}
% -----------------------------------------------------------------------------
We introduce C++ Python bindings that enable the convenient,
efficient, and reliable use of \cgal{} software modules in Python
code. There are different tools that facilitate the creation of C++
Python bindings. We present a short study that compares several tools,
which leads to the method of choice. The binding themselves are
implemented in C++ and their implementation exploits advanced features
in C++, which we explicate. This results with elegant and compact code
that is easy to maintain and extend with additional bindings.
%
We present a system that rapidly generates bindings of desired
objects according to user prescriptions, which enables the convenient
use of any subset of bound object types concurrently.  With our system
it is possible to generate a single library that contains bindings for
instances of different \cgal{} templates, e.g., an instance of the 2D
arrangement class template and an instance of the 2D triangulation class
template, or several libraries when one is insufficient. If several
instances of the same class template must be bound, several
corresponding libraries must be generated, one for each instance.

% -----------------------------------------------------------------------------
\subsection{Outline}
\label{ssec:intro:outline}
% -----------------------------------------------------------------------------
The rest of this paper is organized as follows. A study that compares
several methods for implementing C++ Python bindings is presented
in Section~\ref{sec:comp}.
%
General instructions for rapidly generating bindings with our system 
are presented in Section~\ref{sec:gen}.
%
The description of the bindings generated by our system for various
modules are given in Section~\ref{sec:modules}.
%
A sample of highlighted modern C++ techniques and idioms used in the
binding implementation are described in Section~\ref{sec:code}.
%
The tight coupling between concepts and binding generation is
described in Section~\ref{sec:concepts}.
%
Applications of the new tools, which are geared toward making the CGAL
code more accessible are presented in Section~\ref{sec:app}.

% -----------------------------------------------------------------------------
\subsection{Conventions}
\label{ssec:intro:conventions}
% -----------------------------------------------------------------------------
The paper is packed with code samples written either in the \cmake{},
Python, or C++ languages. We use several conventions to improve
readability.  Code excerpt written in C++ is always numbered, while
short code excerpts written in Python is sometimes preceded with
\pyLstinline{>>>}. The C++ identifier \cppLstinline{bp} is an alias
for the namespace \cppLstinline{boost::python}.
